\SourcePage{236}
\SGAsection{5}{取商与代数结构}
\SGAsection{5.1}{主齐性纤维丛}
\begin{SGAresult}{定义}{5.1.0}
\footnote{编者注：我们引入了编号5.1.0，以便在后文引用。}回顾（III 0.1），具有\Term{group}{群} $H$ 的\Term{right-action}{右作用}的\Term{object}{对象} $X$ 称为在 $H$ 下\Term{formally-principal-homogeneous}{形式主齐性的}\footnote{编者注：也称为“\Term{pseudo-torsor}{伪挠子}”，参见 EGA IV$_4$, 16.5.15。另一方面，更一般的\Term{formally-homogeneous}{形式齐性}\Term{object}{对象}（未必\Term{principal-homogeneous}{主齐性}）的概念，见本讲末尾增补的6.7.1。}，如果由 $(x,h)\mapsto(x,xh)$ 定义的\Term{canonical}{典范}（\Term{functor}{函子}）\Term{morphism}{态射}
\[
X\times H\longrightarrow X\times X
\]
是\Term{isomorphism}{同构}。\Term{equivalence}{等价}地说（参见所引处），对每个 $S\in\operatorname{Ob}\mathcal C$，$X(S)$ 在 $H(S)$ 下\Term{formally-principal-homogeneous}{形式主齐性}，即为空集或在 $H(S)$ 下\Term{principal-homogeneous}{主齐性}。特别地，令 $H$ 通过（右）\Term{translation}{平移}在自身上\Term{action}{作用}，便使 $H$ 成为在自身下\Term{formally-principal-homogeneous}{形式主齐性的}\Term{object}{对象}。
\end{SGAresult}
\begin{SGAresult}{定义}{5.1.1}
具有\Term{group}{群} $H$ 的\Term{action}{作用}的\Term{object}{对象} $X$ 称为\Term{trivial-object-with-operators}{平凡的}，如果它（作为具有 $H$ 的\Term{action}{作用}的\Term{object}{对象}）\Term{isomorphism}{同构}于 $H$，其中 $H$ 通过\Term{translation}{平移}在自身上\Term{action}{作用}。
\end{SGAresult}
\begin{SGAresult}{命题}{5.1.2}
设 $X$ 在 $H$ 下\Term{formally-principal-homogeneous}{形式主齐性}。有一个在 $\Gamma(H)$ 下\Term{principal-homogeneous}{主齐性的}\Term{set}{集合}之间的\Term{isomorphism}{同构}
\[
\Gamma(X)\xrightarrow{\sim}\operatorname{Isom}_{H\text{-obj.}}(H,X).
\]
\end{SGAresult}
\SourcePage{237}
给 $X$ 的每个\Term{section}{截面} $x$，对应从 $H$ 到 $X$ 的\Term{morphism}{态射}，它在\Term{set-sheaf}{集合层}面由 $h\mapsto xh$ 定义。归结到\Term{set}{集合}情形，所述断言立即成立。
\begin{SGAresult}{推论}{5.1.3}
有一个具有 $H$ 的\Term{action}{作用}的\Term{object}{对象}之间的\Term{isomorphism}{同构}
\[
X\xrightarrow{\sim}\underline{\operatorname{Isom}}_{H\text{-obj.}}(H,X).
\]
\end{SGAresult}
\begin{SGAresult}{推论}{5.1.4}
具有\Term{group}{群}的\Term{action}{作用}的\Term{object}{对象}为\Term{trivial-object-with-operators}{平凡的}，当且仅当它\Term{formally-principal-homogeneous}{形式主齐性}并具有一个\Term{section}{截面}。
\end{SGAresult}
\begin{SGAresult}{定义}{5.1.5}
设 $\mathcal C$ 为带有\Term{topology}{拓扑}的\Term{category}{范畴}。具有 $S$-\Term{group}{群} $H$ 的\Term{action}{作用}的 $S$-\Term{object}{对象} $X$ 称为在 $H$ 下的\Term{principal-homogeneous-bundle}{主齐性纤维丛}，如果它\Term{locally-trivial}{局部平凡}，即满足以下\Term{equivalence}{等价}条件：\footnote{编者注：在这种情形，也称 $X$ 为一个 $H$-\Term{torsor}{挠子}。}
\begin{enumerate}
\item[(i)] 使（\Term{functor}{函子}）$X\times_S T$ 在 $H\times_S T$ 下\Term{trivial-object-with-operators}{平凡}的所有 $T\to S$ 所成的\Term{set}{集合}，是 $S$ 的一个\Term{refinement}{加细}。
\item[(ii)] 存在一个\Term{covering-family}{覆盖族}（相对于定义给定\Term{topology}{拓扑}的某个\Term{pretopology}{预拓扑}）$\{S_i\to S\}$，使对每个 $i$，具有 $S_i$-\Term{group-functor}{群函子} $H\times_S S_i$ 的\Term{action}{作用}的 $S_i$-\Term{functor}{函子} $X\times_S S_i$ 都\Term{trivial-object-with-operators}{平凡}（即在 $S_i$ 上有一个\Term{section}{截面}）。
\end{enumerate}
\end{SGAresult}
\begin{SGAresult}{命题}{5.1.6}
设 $\mathcal C$ 为带\Term{topology}{拓扑} $\mathcal T$ 的\Term{category}{范畴}，$(\mathrm M)$ 为 $\mathcal C$ 中满足4.6.3的公理(a)至$(e_{\mathcal T})$的\Term{morphism}{态射}族。设 $H$ 为 $S$-\Term{group}{群}，其\Term{structure-morphism}{结构态射} $H\to S$ 属于 $(\mathrm M)$；$P$ 为具有 $S$-\Term{group}{群} $H$ 的\Term{action}{作用}的 $S$-\Term{object}{对象}。以下条件\Term{equivalence}{等价}：
\begin{enumerate}
\item[(i)] $P$ 是在 $H$ 下的\Term{principal-homogeneous-bundle}{主齐性纤维丛}（定义5.1.5）。
\item[(ii)] $P$ 在 $H$ 下\Term{formally-principal-homogeneous}{形式主齐性}，且\Term{structure-morphism}{结构态射} $P\to S$ 属于 $(\mathrm M)$。
\item[(iii)] 存在属于 $(\mathrm M)$ 的\Term{morphism}{态射} $S'\to S$，使得从 $S$ 到 $S'$ 的\Term{base-extension}{基扩张}使 $P$ 成为\Term{trivial-object-with-operators}{平凡的}，即 $P\times_S S'$ 在 $H\times_S S'$ 下\Term{trivial-object-with-operators}{平凡}。
\item[(iv)] $H$ 在 $P$ 上\Term{free-action}{自由地作用}，且此\Term{action}{作用}为\Term{M-effective}{$(\mathrm M)$-有效的}，而\Term{quotient}{商} $P/H$ \Term{isomorphism}{同构}于 $S$。
\end{enumerate}
\end{SGAresult}
首先注意，(ii)与(iv)\Term{equivalence}{等价}，因为在两种情形下 $P\to S$ 都属于 $(\mathrm M)$，因而\Term{quarrable}{可作基变换}，这保证了\Term{fiber-product}{纤维积} $H\times_S P$ 和 $P\times_S P$ 的\Term{representable}{可表性}。显然(ii)\Term{pushout}{推出}(iii)，因为可取 $S'=P$：$P$ \Term{formally-principal-homogeneous}{形式主齐性}这一假设，蕴含 $P\times_S P$ 在 $H\times_S P$ 下\Term{trivial-object-with-operators}{平凡}（5.1.4），因为它有一个\Term{section}{截面}（\Term{diagonal-morphism}{对角截面}）。显然(iii)\Term{pushout}{推出}(i)，因为由公理$(d_{\mathcal T})$，$\{S'\to S\}$是一个\Term{covering-family}{覆盖族}。因此只须证明(i)\Term{pushout}{推出}(ii)。\Term{sheaf}{层}\Term{morphism}{态射} $P\times_S H\to P\times_S P$ \Term{local-scheme}{局部}为\Term{isomorphism}{同构}，所以为\Term{isomorphism}{同构}（4.5.8）；故 $P$ \Term{formally-principal-homogeneous}{形式主齐性}。\Term{structure-morphism}{结构态射} $P\to S$ \Term{local-scheme}{局部}\Term{isomorphism}{同构}于属于 $(\mathrm M)$ 的\Term{structure-morphism}{结构态射} $H\to S$，故由$(e_{\mathcal T})$，它本身也属于 $(\mathrm M)$。
\SourcePage{238}
(i)与(iv)的\Term{equivalence}{等价}性可作如下推广：
\begin{SGAresult}{命题}{5.1.7}
对 $\mathcal C$ 和 $(\mathrm M)$ 作同样的假设。设 $H$ 为 $S$-\Term{group}{群}，$X$ 为其（从右）\Term{action}{作用}的一个 $S$-\Term{object}{对象}。假设\Term{structure-morphism}{结构态射} $H\to S$ 属于 $(\mathrm M)$。以下条件\Term{equivalence}{等价}：
\begin{enumerate}
\item[(i)] $H$ 在 $X$ 上的\Term{action}{作用}是\Term{free-action}{自由的}且\Term{M-effective}{$(\mathrm M)$-有效的}。
\item[(ii)] 存在一个 $S$-\Term{morphism}{态射} $p:X\to Y$，它与 $H$ 的\Term{action}{作用}在 $X$ 中定义的\Term{equivalence-relation}{等价关系}相容，并且由此\Term{derive}{导出}的 $H\times_S Y$ 在 $Y$ 上的 $X$ 上的\Term{action}{作用}，使 $X$ 成为 $Y$ 上在 $H_Y$ 下的\Term{principal-homogeneous-bundle}{主齐性纤维丛}。
\end{enumerate}
在这些条件下，$p$ 将 $Y$ 等同于\Term{quotient}{商} $X/H$。
\end{SGAresult}
若 $p:X\to Y$ 是与 $H$ 的\Term{action}{作用}相容的\Term{morphism}{态射}，则由公式
\[
H_Y\times_Y X\xrightarrow{\sim}H\times_S X,
\]
\Term{derive}{导出}的 $H\times_S Y$ 在 $Y$ 上的 $X$ 上的\Term{action}{作用}，与 $H$ 的\Term{action}{作用}在 $X$ 中定义相同的\Term{equivalence-relation}{等价关系}。命题由此及前述\Term{equivalence}{等价}性(iv)$\Leftrightarrow$(i)得出。
\begin{SGAresult}{推论}{5.1.7.1}
\footnote{编者注：我们把这个特例作为推论补入；后面的各讲将多次使用它。}设 $\mathcal C$ 为具有\Term{final-object}{终对象}、对\Term{fiber-product}{纤维积}\Term{stable}{稳定}且带有不比\Term{canonical-topology}{典范拓扑}更细的\Term{topology}{拓扑} $\mathcal T$ 的\Term{category}{范畴}。设 $f:G\to H$ 为 $\mathcal C$-\Term{group}{群}的\Term{morphism}{态射}，$K=\operatorname{Ker}(f)$。假设 $f$ 对\Term{topology}{拓扑} $\mathcal T$ 是\Term{covering-morphism}{覆盖态射}。

则 $H$ 表示\Term{quotient-sheaf}{商层} $G/K$，而 $f$ 是一个 $K_H$-\Term{torsor}{挠子}。（也说：“$G$ 是 $H$ 上的一个 $K$-\Term{torsor}{挠子}”。）
\end{SGAresult}
事实上，因 $f$ 是\Term{covering-morphism}{覆盖态射}，它是\Term{universal-effective-epimorphism}{泛有效满态射}（4.4.3）；故由3.3.3.1，$H$ 是 $G$ 关于\Term{equivalence-relation}{等价关系} $\mathscr R(f)=G\times_H G$ 的\Term{quotient}{商}。另一方面，\Term{morphism}{态射} $G\times K\to G\times_H G$，$(g,k)\mapsto(g,gk)$ 是具有 $K_G=G\times_H K_H$ 的\Term{action}{作用}的\Term{object}{对象}之间的\Term{isomorphism}{同构}（其逆由 $(g,g')\mapsto(g,g^{-1}g')$ 给出）。因此，一方面 $\mathscr R(f)$ 就是 $K$ 所定义的\Term{equivalence-relation}{等价关系}；另一方面，因为\Term{morphism}{态射} $f:G\to H$ 是\Term{covering-morphism}{覆盖态射}，由5.1.6(ii)（或直接由定义5.1.5(ii)），$f$ 是一个 $K_H$-\Term{torsor}{挠子}。

现在可在关于\Term{group}{群}\Term{action}{作用}的取\Term{quotient}{商}情形，将\Term{theorem}{定理}4.6.6进一步明确如下：
\begin{SGAresult}{命题}{5.1.8}
在5.1.7的假设下，记 $F_0$ 为 $S$ 上如下定义的\Term{functor}{函子}：对每个 $S'\to S$，$F_0(S')$ 是 $X\times_S S'$ 的所有\Term{representable}{可表}子 $S'$-\Term{functor}{函子} $Z$ 所成的\Term{set}{集合}，其中 $Z$ 在 $H\times_S S'$ 下\Term{stable}{稳定}，并在\Term{induced-action}{诱导作用}下成为这一 $S'$-\Term{group}{群}下的\Term{principal-homogeneous-bundle}{主齐性纤维丛}（3.2.2）。
\begin{enumerate}
\item[(i)] 以下条件\Term{equivalence}{等价}：
\begin{enumerate}
\item[a)] $H$ 在 $X$ 中的\Term{action}{作用}为\Term{M-effective}{$(\mathrm M)$-有效的}且\Term{free-action}{自由的}\footnote{编者注：我们补入了“且自由”。}。
\item[b)] $F_0$ \Term{representable}{可表}。
\end{enumerate}
在这些条件下，$F_0=X/H$。
\SourcePage{239}
\item[(ii)] 设 $(\mathrm N)$ 为对\Term{base-change}{基变换}\Term{stable}{稳定}的一族\Term{morphism}{态射}，并且对每个\Term{covering-family}{覆盖族} $\{S'_i\to S'\}$ 及属于 $(\mathrm N)$ 的任一\Term{morphism}{态射}族 $\{T_i\to S'_i\}$，$T_i$ 上相对于 $\{S'_i\to S'\}$ 的每个\Term{descent-datum}{下降资料}都是\Term{effective-descent-datum}{有效的}。假设\Term{morphism}{态射} $X\times_S H\to X\times_S X$ 属于 $(\mathrm N)$，且 $X$ \Term{quarrable}{可作基变换}。则与 $X\times_S X$ 的\Term{subobject}{子对象} $X\times_S H$ 对应的 $\operatorname{Hom}(X,F_0)$ 中的元素 $p$，将 $F_0$ 等同于\Term{quotient-sheaf}{商层} $X/H$。
\end{enumerate}
\end{SGAresult}
\SGAsection{5.2}{群结构与取商}
本小节研究\Term{group}{群}对其\Term{subgroup}{子群}的\Term{quotient}{商} $G/H$ 上可以赋予的\Term{algebraic-structure}{代数结构}。我们先在 $\mathcal C$ 上关于任意\Term{topology}{拓扑}的\Term{sheaf}{层}\Term{category}{范畴}中讨论。取\Term{canonical-topology}{典范拓扑}并利用4.5.12，将得到 $\mathcal C$ 中关于\Term{universal-effective-quotient}{泛有效取商}的结果。
\begin{SGAresult}{命题}{5.2.1}
设 $u:H\to G$ 为\Term{group-sheaf}{群层}的\Term{monomorphism}{单态射}。\Term{quotient-sheaf}{商层} $G/H$ 上存在唯一的具有\Term{group}{群} $G$ 的\Term{action}{作用}的\Term{object}{对象}结构，使\Term{canonical}{典范}\Term{morphism}{态射}
\[
p:G\longrightarrow G/H
\]
是具有\Term{group}{群} $G$ 的\Term{action}{作用}的\Term{object}{对象}之间的\Term{morphism}{态射}。这一结构关于二元组 $(G,H)$ 具有\Term{functorial}{函子性}：若有\Term{commutative-diagram}{交换图}
\[
\begin{tikzcd}
H\arrow[r]\arrow[d]&G\arrow[d,"f"]\\
H'\arrow[r]&G',
\end{tikzcd}
\]
则\Term{morphism}{态射} $\bar f:G/H\to G'/H'$（3.2.3）与\Term{action}{作用}\Term{group}{群}上的\Term{morphism}{态射} $f$ 相容。
\end{SGAresult}
事实上，\Term{sheaf}{层} $G/H$ 是\Term{presheaf}{预层}
\[
i(G)/i(H):S\longmapsto G(S)/H(S)
\]
的\Term{associated-sheaf}{伴随层}；因为\Term{functor}{函子} $a$ \Term{left-exact}{左正合}，它将具有\Term{group}{群}\Term{action}{作用}的\Term{object}{对象}变为具有\Term{group}{群}\Term{action}{作用}的\Term{object}{对象}。\Term{presheaf}{预层} $i(G)/i(H)$ 具有\Term{group}{群} $i(G)$ 的\Term{action}{作用}，故 $G/H=a(i(G)/i(H))$ 具有\Term{group}{群} $a(i(G))=G$ 的\Term{action}{作用}。这一结构显然具有所述全部性质。
\begin{SGAresult}{推论}{5.2.2}
设 $u:H\to G$ 为 $\mathcal C$-\Term{group}{群}的\Term{monomorphism}{单态射}。假设 $H$ 在 $G$ 上的\Term{action}{作用}为\Term{universal-effective-action}{泛有效的}。\Term{quotient-object}{商对象} $G/H\in\operatorname{Ob}\mathcal C$ 上存在唯一的具有\Term{group}{群} $G$ 的\Term{action}{作用}的\Term{object}{对象}结构，使 $p:G\to G/H$ 是带\Term{action}{作用}\Term{object}{对象}的\Term{morphism}{态射}。这一结构关于二元组 $(H,G)$（$H$ 在 $G$ 中的\Term{action}{作用}为\Term{universal-effective-action}{泛有效的}）在上述意义下具有\Term{functorial}{函子性}。
\end{SGAresult}
\SourcePage{240}
\begin{SGAresult}{命题}{5.2.3}
设 $u:H\to G$ 为\Term{group-sheaf}{群层}的\Term{monomorphism}{单态射}，将 $H$ 等同于 $G$ 的一个\Term{invariant}{不变}\Term{subgroup-sheaf}{子群层}。\Term{quotient-sheaf}{商层} $G/H$ 上存在唯一的\Term{group-sheaf}{群层}结构，使\Term{canonical}{典范}\Term{morphism}{态射} $p:G\to G/H$ 为\Term{group-morphism}{群态射}。这一结构关于二元组 $(H,G)$（$H$ \Term{invariant}{不变}）具有\Term{functorial}{函子性}。
\end{SGAresult}
证明与5.2.1的证明相似。
\begin{SGAresult}{推论}{5.2.4}
设 $u:H\to G$ 为 $\mathcal C$-\Term{group}{群}的\Term{monomorphism}{单态射}，将 $H$ 等同于 $G$ 的一个\Term{invariant}{不变}\Term{subgroup}{子群}。假设 $H$ 在 $G$ 上的\Term{action}{作用}为\Term{universal-effective-action}{泛有效的}。\Term{quotient-object}{商对象} $G/H\in\operatorname{Ob}\mathcal C$ 上存在唯一的\Term{group}{群}结构，使\Term{canonical}{典范}\Term{morphism}{态射} $G\to G/H$ 为\Term{group-morphism}{群态射}。这一结构关于二元组 $(H,G)$（$H$ \Term{invariant}{不变}，且其\Term{action}{作用}为\Term{universal-effective-action}{泛有效的}）具有\Term{functorial}{函子性}。
\end{SGAresult}
$G/H$ 的\Term{group}{群}结构可作如下更直观的刻画：
\begin{SGAresult}{命题}{5.2.5}
在5.2.4的条件下，设 $K$ 为 $\mathcal C$-\Term{group}{群}，$f:G\to K$ 为\Term{morphism}{态射}。以下条件\Term{equivalence}{等价}：
\begin{enumerate}
\item[(i)] $f$ 是与 $H$ 所定义的\Term{equivalence-relation}{等价关系}相容的\Term{group-morphism}{群态射}。
\item[(ii)] $f$ 是诱导\Term{trivial-morphism}{平凡态射} $H\to K$ 的\Term{group-morphism}{群态射}。
\item[(iii)] $f$ \Term{resolution}{分解}为一个\Term{group-morphism}{群态射} $G/H\to K$。
\end{enumerate}
特别地，有关于\Term{group}{群} $K$ 具有\Term{functorial}{函子性}的\Term{isomorphism}{同构}
\[
\operatorname{Hom}_{\mathcal C\text{-gr.}}(G/H,K)\xrightarrow{\sim}
\{f\in\operatorname{Hom}_{\mathcal C\text{-gr.}}(G,K)\mid f\circ u=e\}.
\]
\end{SGAresult}
(i)与(ii)的\Term{equivalence}{等价}性在\Term{set-sheaf}{集合层}面证明。显然(iii)$\Rightarrow$(ii)。(iii)与(ii)的\Term{equivalence}{等价}性由公式
\[
\operatorname{Hom}(G/H,K)\simeq\operatorname{Hom}(i(G)/i(H),K)
\]
及 $G/H$ 的\Term{group}{群}结构的定义得出。
\begin{SGAresult}{注记}{5.2.6}
在上述情形，若 $f$ 的\Term{kernel}{核}恰为 $H$，则\Term{resolution}{分解} $f$ 的\Term{morphism}{态射} $G/H\to K$ 为\Term{monomorphism}{单态射}。这由3.3.4立即得出。
\end{SGAresult}
对于\Term{group-sheaf}{群层}，4.4.11可进一步明确为：
\begin{SGAresult}{命题}{5.2.7}
设 $G$ 为\Term{group-sheaf}{群层}，$H$ 为其\Term{invariant}{不变}\Term{subgroup-sheaf}{子群层}。对 $G$ 的每个包含 $H$ 的\Term{subgroup-sheaf}{子群层} $K$，将\Term{quotient-group}{商群} $K/H$ 记为 $K'$，并将它看作 $G'=G/H$ 的一个\Term{subgroup}{子群}。

有 $K=K'\times_{G'}G$，映射 $K\mapsto K/H$ 与 $K'\mapsto K'\times_{G'}G$ 给出 $G$ 的包含 $H$ 的\Term{subgroup-sheaf}{子群层}的\Term{set}{集合}与 $G'$ 的\Term{subgroup-sheaf}{子群层}的\Term{set}{集合}之间的\Term{bijection}{双射}。在此对应下，$G$ 与 $G'$ 的\Term{invariant}{不变}\Term{subgroup-sheaf}{子群层}相互对应。
\end{SGAresult}
\SourcePage{241}
第一部分易由4.4.11与3.2.4得出。只须说明，$K$ 在 $G$ 中\Term{invariant}{不变}当且仅当 $K'$ 在 $G'$ 中\Term{invariant}{不变}。若 $K$ 在 $G$ 中\Term{invariant}{不变}，则\Term{presheaf}{预层} $i(K)/i(H)$ 在 $i(G)/i(H)$ 中\Term{invariant}{不变}；由惯常的论证，其\Term{associated-sheaf}{伴随层}也如此。反过来，若 $K'$ 在 $G'$ 中\Term{invariant}{不变}，则\Term{fiber-product}{纤维积} $K\times_{G'}G$ 在 $G$ 中\Term{invariant}{不变}，这是立即可见的。

现在若 $L$ 为 $G$ 的任意\Term{subgroup-sheaf}{子群层}，记 $\overline L$ 为 $L$ 关于 $H$ 所定义的\Term{equivalence-relation}{等价关系}的\Term{saturation}{饱和化}，也记作 $\overline L=L\cdot H$。
\begin{SGAresult}{命题}{5.2.8}
在上述条件下，$L\cdot H$ 是 $G$ 的一个包含 $H$ 的\Term{subgroup-sheaf}{子群层}，$L$ 在 $G/H$ 中的像等同于
\[
(L\cdot H)/H\simeq L/(H\cap L).
\]
\end{SGAresult}
事实上，记 $L'$ 为 $L$ 在 $G/H$ 中的像\Term{sheaf}{层}。它是 $G/H$ 的一个\Term{subgroup-sheaf}{子群层}，在上一命题的对应中对应于 $L\cdot H$。因为\Term{morphism}{态射} $L\to L'$ 是\Term{covering-morphism}{覆盖态射}，因而是\Term{sheaf}{层}的\Term{universal-effective-epimorphism}{泛有效满态射}，由4.4.9，$L'$ 等同于 $L$ 对 $L\to L'$ 的\Term{kernel}{核}的\Term{quotient}{商}，而此\Term{kernel}{核}显然就是 $H\cap L$。

最后考虑如下情形：给定\Term{group-sheaf}{群层} $G$、其\Term{subgroup-sheaf}{子群层} $K$，及 $K$ 中\Term{invariant}{不变}的\Term{subgroup-sheaf}{子群层} $H$。先定义\Term{group-sheaf}{群层} $H\backslash K\;(=K/H)$ 在 $G/H$ 上的\Term{right-action}{右作用}。\Term{group}{群} $K$ 通过\Term{right-translation}{右平移}在 $G$ 上\Term{action}{作用}。因为 $H$ 在 $K$ 中\Term{invariant}{不变}，这一\Term{action}{作用}与 $H$ 的\Term{action}{作用}所定义的\Term{equivalence-relation}{等价关系}相容，因而定义了 $K$ 在 $G/H$ 上的\Term{action}{作用}，即从 $K$ 的\Term{opposite-group}{反群} $K^{\circ}$ 到 $\underline{\operatorname{Aut}}(G/H)$ 的\Term{morphism}{态射}。后者是一个\Term{sheaf}{层}（4.5.13），且这一\Term{morphism}{态射}在 $H$ 上\Term{trivial-morphism}{平凡}，故它经 $K/H$ \Term{resolution}{分解}，并定义所求的\Term{action}{作用}。因为 $G$ 在自身上的\Term{right-action}{右作用}与\Term{left-action}{左作用}可交换，$G$ 和 $K/H$ 在 $G/H$ 上的\Term{action}{作用}也可交换。
\begin{SGAresult}{命题}{5.2.9}
在上述条件下，$K/H$ 在 $G/H$ 中\Term{free-action}{自由地}（从右）\Term{action}{作用}，并且有具有\Term{group}{群} $G$ 的\Term{action}{作用}的\Term{sheaf}{层}之间的\Term{canonical}{典范}\Term{isomorphism}{同构}
\[
(G/H)/(K/H)\simeq G/K.
\]
若 $K$ 在 $G$ 中\Term{invariant}{不变}，此时 $K/H$ 在 $G/H$ 中也\Term{invariant}{不变}（5.2.7），则这一\Term{isomorphism}{同构}保持两边的\Term{group}{群}结构。
\end{SGAresult}
有\Term{presheaf}{预层}的\Term{isomorphism}{同构}
\[
\bigl(i(G)/i(H)\bigr)\big/\bigl(i(K)/i(H)\bigr)
\xleftarrow{\sim}i(G)/i(K),
\]
它保持具有\Term{group}{群} $i(G)$ 的\Term{action}{作用}的\Term{object}{对象}结构。对此关系应用\Term{functor}{函子} $a$，即得所述结果。
\begin{SGAresult}{推论}{5.2.10}
设 $G$ 为 $\mathcal C$-\Term{group}{群}，$K$ 为 $G$ 的一个子 $\mathcal C$-\Term{group}{群}，$H$ 为 $K$ 的一个\Term{invariant}{不变}子 $\mathcal C$-\Term{group}{群}。设 $(\mathrm M)$ 为 $\mathcal C$ 中满足公理(a)至$(e_{\mathcal T})$的一族\Term{morphism}{态射}。假设 $H$ 在 $G$（相应地，$K$）上的\Term{right-action}{右作用}为\Term{M-effective}{$(\mathrm M)$-有效的}。则 $K/H$ 自然地从右\Term{free-action}{自由作用}于 $G/H$；这一\Term{action}{作用}与 $G$ 的\Term{action}{作用}可交换。以下条件\Term{equivalence}{等价}：
\SourcePage{242}
\begin{enumerate}
\item[(i)] $K$ 在 $G$ 上的\Term{action}{作用}为\Term{M-effective}{$(\mathrm M)$-有效的}。
\item[(ii)] $K/H$ 在 $G/H$ 上的\Term{action}{作用}为\Term{M-effective}{$(\mathrm M)$-有效的}。
\end{enumerate}
在这些条件下，有具有\Term{group}{群} $G$ 的\Term{action}{作用}的\Term{object}{对象}\footnote{编者注：指 $\mathcal C$ 中的\Term{object}{对象}。}之间的\Term{isomorphism}{同构}
\[
(G/H)/(K/H)\simeq G/K.
\]
\end{SGAresult}
\SGAsection{5.3}{有效性判据的应用：Noether定理}
\footnote{编者注：实际上，5.2.8至5.2.10中建立的\Term{isomorphism}{同构}更应称为“Noether\Term{isomorphism}{同构}\Term{theorem}{定理}”。}设 $\mathcal C$、$\mathcal T$ 与 $(\mathrm M)$ 如通常情形。设 $(\mathrm N)$ 为满足4.7的公理(a)和$(f_{\mathrm M})$的一族\Term{morphism}{态射}。结合5.2.7与4.7.2，得到：
\begin{SGAresult}{命题}{5.3.1}
设 $G$ 为 $\mathcal C$-\Term{group}{群}，$H$ 为 $G$ 的一个\Term{invariant}{不变}子 $\mathcal C$-\Term{group}{群}，其在 $G$ 中的\Term{action}{作用}为\Term{M-effective}{$(\mathrm M)$-有效的}。

对 $G$ 的每个包含 $H$ 且\Term{morphism}{态射} $K\to G$ 属于 $(\mathrm N)$ 的子 $\mathcal C$-\Term{group}{群} $K$，$H$ 在 $K$ 中的\Term{action}{作用}为\Term{M-effective}{$(\mathrm M)$-有效的}，并且\Term{quotient}{商} $K/H=K'$ 是 $G/H=G'$ 的一个子 $\mathcal C$-\Term{group}{群}，其\Term{morphism}{态射} $K'\to G'$ 属于 $(\mathrm N)$。

映射 $K\mapsto K'$ 给出 $G$ 的包含 $H$ 且 $K\to G$ 属于 $(\mathrm N)$ 的子 $\mathcal C$-\Term{group}{群} $K$ 的\Term{set}{集合}，与 $G'$ 的 $K'\to G'$ 属于 $(\mathrm N)$ 的子 $\mathcal C$-\Term{group}{群} $K'$ 的\Term{set}{集合}之间的\Term{bijection}{双射}。其逆映射为 $K'\mapsto K\times_{G'}G$。在此对应下，$G$ 与 $G'$ 的\Term{invariant}{不变}\Term{subgroup}{子群}相互对应。
\end{SGAresult}
\begin{SGAresult}{推论}{5.3.2}
若 $H\to G$ 属于 $(\mathrm N)$，则 $\mathcal C$ 具有\Term{final-object}{终对象} $e$，且\Term{unit-section}{单位截面} $e\to G/H$ 属于 $(\mathrm N)$。\footnote{编者注：这是将命题应用于 $K=H$ 所得的结果。}
\end{SGAresult}
\SGAsection{6}{概形范畴中的拓扑}
\SGAsection{6.1}{Zariski拓扑}
这是由如下\Term{pretopology}{预拓扑}生成的\Term{topology}{拓扑}：\Term{morphism}{态射}族 $\{S_i\to S\}$ 为\Term{covering-family}{覆盖族}，如果每个\Term{morphism}{态射}都是\Term{open-immersion}{开浸入}，且 $S_i$ 的像之并为整个 $S$。记此\Term{topology}{拓扑}为(Zar)。
\begin{SGAresult}{定义}{6.1.1}
\Term{zariski-topology}{Zariski拓扑}的\Term{sheaf}{层}也称为\Term{local-nature-functor}{局部性函子}：它是从$(\mathrm{Sch})$到$(\mathrm{Ens})$的\Term{contravariant-functor}{反变函子}，使对每个\Term{scheme}{概形} $S$ 及 $S$ 的每个由\Term{open-subset}{开集} $S_i$ 构成的\Term{covering}{覆盖}，都有\Term{exact}{正合}图式
\[
F(S)\longrightarrow\prod_i F(S_i)\rightrightarrows\prod_{i,j}F(S_i\cap S_j).
\]
\end{SGAresult}
特别地，\Term{local-nature-functor}{局部性函子}将\Term{direct-sum}{直和}变为\Term{product}{积}。因为每个\Term{representable}{可表}\Term{functor}{函子}都是\Term{sheaf}{层}，这一\Term{topology}{拓扑}不比\Term{canonical-topology}{典范拓扑}更细。
\SourcePage{243}
术语上，凡不加说明而说“\Term{local-scheme}{局部}的”“\Term{local-scheme}{局部}地”，均指相对于\Term{zariski-topology}{Zariski拓扑}，即通常的含义。
\SGAsection{6.2}{一种构造拓扑的方法}
\begin{SGAresult}{命题}{6.2.1}
设 $\mathcal C$ 为\Term{category}{范畴}，$\mathcal C'$ 为其\Term{full-subcategory}{满子范畴}，$P$ 为 $\mathcal C$ 中具有相同目标的\Term{morphism}{态射}族所成的一个\Term{set}{集合}，它对\Term{base-change}{基变换}及复合\Term{stable}{稳定}（即满足4.2.5的公理(P1)与(P2)）；$P'$ 为 $\mathcal C'$ 中\Term{morphism}{态射}族所成的一个\Term{set}{集合}，包含仅由一个恒等\Term{isomorphism}{同构}构成的各族。在 $\mathcal C$ 上赋予由 $P$ 和 $P'$ 生成的\Term{topology}{拓扑}（参见4.2.5.0），并假设下列三个条件成立：
\begin{enumerate}
\item[(a)] 若 $\{S_i\to S\}\in P'$（故 $S_i,S\in\operatorname{Ob}\mathcal C'$），且 $T\to S$ 为 $\mathcal C'$ 中的\Term{morphism}{态射}，则（$\mathcal C$ 中的）\Term{fiber-product}{纤维积} $S_i\times_S T$ 存在，且族 $\{S_i\times_S T\to T\}$ 属于 $P'$（故 $S_i\times_S T\in\operatorname{Ob}\mathcal C'$）。（注意：这一条件蕴含 $P'$ 对 $\mathcal C'$ 中的\Term{base-change}{基变换}\Term{stable}{稳定}，但并不与之\Term{equivalence}{等价}，因为它还要求从 $\mathcal C'$ 到 $\mathcal C$ 的包含\Term{functor}{函子}与某些\Term{fiber-product}{纤维积}可交换。）
\item[(b)] 对每个 $S\in\operatorname{Ob}\mathcal C$，存在 $\{S_i\to S\}\in P$，使每个 $S_i\in\operatorname{Ob}\mathcal C'$。
\item[(c)] 在如下情形中：
\[
\begin{tikzcd}
&S_{ijk}\arrow[d,"(P')"]\\
S_i\arrow[d,"(P')"']&S_{ij}\arrow[l,"(P)"']\\
S&
\end{tikzcd}
\]
其中 $S,S_i,S_{ij},S_{ijk}\in\operatorname{Ob}\mathcal C'$；$\{S_i\to S\}\in P'$；对每个 $i$，$\{S_{ij}\to S_i\}\in P$；对每个 $ij$，$\{S_{ijk}\to S_{ij}\}\in P'$，则存在一个族 $\{T_n\to S\}\in P'$，并且对每个 $n$，存在一个多重指标 $(ijk)$ 及\Term{commutative-diagram}{交换图}
\[
\begin{tikzcd}
S_{ijk}\arrow[dr]&T_n\arrow[l]\arrow[d]\\
&S.
\end{tikzcd}
\]
\end{enumerate}
那么，$S\in\operatorname{Ob}\mathcal C$ 的一个\Term{sieve}{筛} $R$ 为\Term{covering-sieve}{覆盖筛}，当且仅当存在复合族
\[
\begin{tikzcd}
R\arrow[d]&S_{ij}\arrow[l,dashed]\arrow[d,"(P')"]\\
S&S_i\arrow[l,"(P)"']
\end{tikzcd}
\]
\SourcePage{244}
其中 $S_i,S_{ij}\in\operatorname{Ob}\mathcal C'$，$\{S_i\to S\}\in P$，对每个 $i$ 有 $\{S_{ij}\to S_i\}\in P'$，并且所得\Term{morphism}{态射} $S_{ij}\to S$ 经 $R$ \Term{resolution}{分解}（换言之，这一复合族生成的\Term{sieve}{筛}包含在 $R$ 中）。
\end{SGAresult}
\begin{proof}
因为属于 $P$ 与 $P'$ 的族都是\Term{covering-family}{覆盖族}，这些族的复合族也为\Term{covering-family}{覆盖族}（C2），所以具有所述\Term{form}{形式}的\Term{sieve}{筛}包含一个\Term{covering-sieve}{覆盖筛}，从而也是\Term{covering-sieve}{覆盖筛}。

反过来，只须说明具有此\Term{form}{形式}的\Term{sieve}{筛}确实构成一个\Term{topology}{拓扑}，即只须验证4.2.1的公理(T1)至(T4)。

\textit{公理(T4)。}设 $S\in\operatorname{Ob}\mathcal C$。由(b)，存在族 $\{S_i\to S\}\in P$，且 $S_i\in\operatorname{Ob}\mathcal C'$。由假设，族 $\{S_i\xrightarrow{\mathrm{id}_{S_i}}S_i\}$ 属于 $P'$。因此，$S$ 的\Term{sieve}{筛} $S$ 具有所需\Term{form}{形式}：
\[
\begin{tikzcd}
S_i\arrow[r]\arrow[d,"(P')\;\mathrm{id}"']&S\arrow[d,"\mathrm{id}_S"]\\
S_i\arrow[r,"(P)"']&S.
\end{tikzcd}
\]
\textit{公理(T3)。}显然。

\textit{公理(T2)。}设 $R$ 为 $S$ 的一个具有所需\Term{form}{形式}的\Term{sieve}{筛}，$C$ 为 $S$ 的一个\Term{sieve}{筛}\footnote{编者注：我们在此修正了原文。}，使对每个 $T\in\operatorname{Ob}\mathcal C$ 及每个经 $R$ \Term{resolution}{分解}的\Term{morphism}{态射} $T\to S$，$T$ 的\Term{sieve}{筛} $C\times_S T$ 都具有所需\Term{form}{形式}。那么，由于 $S_{ij}\to S$ 经 $R$ \Term{resolution}{分解}，$S_{ij}$ 的\Term{sieve}{筛} $C_{ij}=C\times_S S_{ij}$
\[
\begin{tikzcd}[column sep=large]
&S_{ij}\arrow[d,"(P')"]&C_{ij}\arrow[l,hook']\arrow[dd]\\
&S_i\arrow[d,"(P)"]&\\
R\arrow[r,hook]&S&C\arrow[l,hook']
\arrow[from=1-2,to=3-1]
\end{tikzcd}
\]
具有所需\Term{form}{形式}；故对每个 $ij$，有如下\Term{form}{形式}的图式：
\[
\begin{tikzcd}
S_{ijkl}\arrow[d,"(P')"']\arrow[ddr] &\\
S_{ijk}\arrow[d,"(P)"']&\\
S_{ij}&C_{ij}\arrow[l,hook'].
\end{tikzcd}
\]
\SourcePage{245}
我们因此证明了，存在一个复合族
\[
S_{ijkl}\xrightarrow{(P')}S_{ijk}\xrightarrow{(P)}S_{ij}
\xrightarrow{(P')}S_i\xrightarrow{(P)}S,
\]
它属于 $P\circ P'\circ P\circ P'$，经 $C$ \Term{resolution}{分解}，并且除 $S$ 之外的所有\Term{object}{对象}都在 $\mathcal C'$ 中。对每个族 $\{S_{ijkl}\to S_i\}$ 应用条件(c)，可知对每个 $i$ 存在一个族 $\{T_{in}\to S_i\}\in P'$，使 $T_{in}\to S$ 经某个 $S_{ijkl}$ \Term{resolution}{分解}，从而经 $C$ \Term{resolution}{分解}：
\[
\begin{tikzcd}
T_{in}\arrow[r]\arrow[d,"(P')"']&S_{ijkl}\arrow[dd]\\
S_i\arrow[d,"(P)"']&\\
S&C\arrow[l,hook'].
\end{tikzcd}
\]
故 $S$ 的\Term{sieve}{筛} $C$ 具有所需\Term{form}{形式}，验证完毕。

\textit{公理(T1)。}设 $R$ 为 $S$ 的一个具有给定\Term{form}{形式}的\Term{sieve}{筛}，$T\to S$ 为 $\mathcal C$ 中的\Term{morphism}{态射}。证明 $T$ 的\Term{sieve}{筛} $R\times_S T$ 具有所需\Term{form}{形式}。
\[
\begin{tikzcd}[column sep=large]
&S_{ij}\arrow[d,"(P')"]&&U_{ikj}\arrow[ll]\arrow[d,"(P')"]\\
&S_i\arrow[d,"(P)"]&T_i\arrow[l]\arrow[d,"(P)"]&U_{ik}\arrow[l,"(P)"']\arrow[dl,"(P)"]\\
R\arrow[r,hook]&S&T\arrow[l]&
\arrow[from=1-2,to=3-1]
\end{tikzcd}
\]
令 $T_i=S_i\times_S T$。族 $\{T_i\to T\}$ 属于 $P$（由(P1)）。应用(b)，构造 $\{U_{ik}\to T_i\}\in P$，其中 $U_{ik}\in\operatorname{Ob}\mathcal C'$。由假设（$P$ 的条件(P2)），有 $\{U_{ik}\to T\}\in P$。由(a)，$U_{ik}\times_{S_i}S_{ij}=U_{ikj}$ 是 $\mathcal C'$ 的\Term{object}{对象}，且对每个 $ik$，$\{U_{ikj}\to U_{ik}\}\in P'$。于是下列\Term{commutative-diagram}{交换图}
\[
\begin{tikzcd}
R\arrow[d]&&U_{ikj}\arrow[ll]\arrow[d,"(P')"]\\
S&T\arrow[l]&U_{ik}\arrow[l,"(P)"']
\end{tikzcd}
\]
表明\Term{morphism}{态射} $U_{ijk}\to T$ 经 $T$ 的\Term{sieve}{筛} $T\times_S R$ \Term{resolution}{分解}，所以此\Term{sieve}{筛}具有所需\Term{form}{形式}。证明完毕。
\end{proof}
\begin{SGAresult}{推论}{6.2.2}
若 $S\in\operatorname{Ob}\mathcal C'$，且 $R$ 为 $S$ 的一个\Term{sieve}{筛}，则 $R$ 为\Term{covering-sieve}{覆盖筛}当且仅当存在一个经 $R$ \Term{resolution}{分解}的族 $\{T_i\to S\}\in P'$。
\end{SGAresult}
\SourcePage{246}
事实上，这样的\Term{sieve}{筛}为\Term{covering-sieve}{覆盖筛}。另一方面，取族 $\{S_i\to S\}$ 仅由 $S$ 的恒等\Term{isomorphism}{同构}构成，应用(c)，即可从命题\Term{pushout}{推出}\Term{covering-sieve}{覆盖筛}具有所述\Term{form}{形式}。
\begin{SGAresult}{推论}{6.2.3}
$\mathcal C$ 上的\Term{presheaf}{预层} $F$ 为\Term{separated-presheaf}{分离的}（相应地，为\Term{sheaf}{层}），当且仅当在以下两种情形中，\Term{morphism}{态射}
\[
F(S)\longrightarrow\prod_i F(S_i)
\]
为\Term{injectivity}{单射}（相应地，图式
\[
F(S)\longrightarrow\prod_i F(S_i)\rightrightarrows\prod_{i,j}F(S_i\times_S S_j)
\]
\Term{exact}{正合}）：
\begin{enumerate}
\item[(i)] $\{S_i\to S\}\in P$；
\item[(ii)] $S,S_i\in\operatorname{Ob}\mathcal C'$；$\{S_i\to S\}\in P'$。
\end{enumerate}
\end{SGAresult}
这些条件是必要的，因为所论各族都是\Term{covering-family}{覆盖族}。若 $C$ 为\Term{morphism}{态射}族 $\{S_{ij}\xrightarrow{(P')}S_i\xrightarrow{(P)}S\}$ 的像\Term{sieve}{筛}，则简单的图追踪表明，推论的条件蕴含 $\operatorname{Hom}(S,F)\xrightarrow{g}\operatorname{Hom}(C,F)$ 为\Term{injectivity}{单射}（相应地，为\Term{bijection}{双射}）。而 $S$ 的每个\Term{refinement}{加细} $R$ 都包含一个上述类型的\Term{sieve}{筛} $C$，并且有\Term{commutative-diagram}{交换图}
\[
\begin{tikzcd}[column sep=large]
\operatorname{Hom}(S,F)\arrow[r,"f"]\arrow[dr,"g"']&\operatorname{Hom}(R,F)\arrow[d,"h"]\\
&\operatorname{Hom}(C,F).
\end{tikzcd}
\]
已知 $g$ 为\Term{injectivity}{单射}，故 $f$ 亦然。因此 $F$ \Term{separated-presheaf}{分离}。但 $R$ 是 $C$ 的一个\Term{refinement}{加细}，故 $h$ 也为\Term{injectivity}{单射}。若 $g$ 为\Term{bijection}{双射}，则 $f$ 也为\Term{bijection}{双射}，所以 $F$ 是\Term{sheaf}{层}。
\begin{SGAresult}{注记}{6.2.4}
上一推论不能由4.3.5得出，因为 $P'$ 对\Term{base-extension}{基扩张}并不\Term{stable}{稳定}。
\end{SGAresult}
\begin{SGAresult}{注记}{6.2.5}
特别地，在如下情形下条件(c)成立：
\begin{enumerate}
\item[(i)] $P'$ 对复合\Term{stable}{稳定}。
\item[(ii)] 若 $\{S_j\to S\}$ 为 $\mathcal C'$ 中的\Term{morphism}{态射}族且属于 $P$，则它有一个属于 $P'$ 的子族。
\end{enumerate}
\end{SGAresult}
\SGAsection{6.3}{在概形范畴中的应用}
取 $\mathcal C$ 为\Term{scheme}{概形}\Term{category}{范畴}，$\mathcal C'$ 为由\Term{affine}{仿射}\Term{scheme}{概形}构成的\Term{full-subcategory}{满子范畴}，$P$ 为\Term{open-immersion}{开浸入}的所有\Term{surjective-family}{满射族}所成的\Term{set}{集合}。我们将考虑若干\Term{set}{集合} $P'$：
\begin{itemize}
\item $P'_1$：由\Term{flat}{平坦}\Term{morphism}{态射}构成的\Term{finite}{有限}\Term{surjective-family}{满射族}。
\SourcePage{247}
\item $P'_2$：由\Term{finite-presentation}{有限表示}且\Term{quasi-finite}{拟有限}的\Term{flat}{平坦}\Term{morphism}{态射}构成的\Term{finite}{有限}\Term{surjective-family}{满射族}。\footnote{编者注：设 $P''_2$ 为由 $\mathcal C'$ 中\Term{finite-presentation}{有限表示}的\Term{flat}{平坦}\Term{morphism}{态射}构成的\Term{finite}{有限}\Term{surjective-family}{满射族}所成的\Term{set}{集合}。由下面的命题6.3.1，$P$ 与 $P'_2$ 生成的\Term{topology}{拓扑} $\mathcal T_2$，等于 $P$ 与 $P''_2$ 生成的\Term{topology}{拓扑}。这是 EGA IV$_4$, \S17.16 关于\Term{quasi-section}{拟截面}的结果的推论，参见6.3.1的证明。此处引述 EGA IV$_4$, 17.16.2 的如下特例：设 $S$ \Term{affine}{仿射}，$f:X\to S$ 为\Term{locally-finite-presentation}{局部有限表示}的\Term{flat}{平坦}\Term{surjectivity}{满射}\Term{morphism}{态射}，则存在\Term{faithfully-flat}{忠实平坦}、\Term{finite-presentation}{有限表示}且\Term{quasi-finite}{拟有限}的\Term{morphism}{态射} $S'\to S$，其中 $S'$ \Term{affine}{仿射}，以及一个 $S$-\Term{morphism}{态射} $S'\to X$。}
\item $P'_3$：由\Term{etale}{平展}\Term{morphism}{态射}构成的\Term{finite}{有限}\Term{surjective-family}{满射族}。
\item $P'_4$：由\Term{finite}{有限}\Term{etale}{平展}\Term{morphism}{态射}构成的\Term{finite}{有限}\Term{surjective-family}{满射族}。
\end{itemize}
除 $P'_4$ 外，每个\Term{set}{集合} $P'_i$ 都满足命题6.2.1的条件（(c)由6.2.5得出，因为\Term{affine}{仿射}\Term{scheme}{概形}是\Term{quasi-compact}{拟紧}的，$\mathcal C'$ 中每个属于 $P$ 的\Term{morphism}{态射}族都包含一个也属于 $P$ 的\Term{finite}{有限}子族，因而对 $i=1,2,3$，此子族属于 $P'_i$）。由 $P$ 与 $P'_i$ 生成的\Term{topology}{拓扑} $\mathcal T_i$ 的记号与名称如下\footnotemark[\value{footnote}]：
\begin{center}
\begin{tabular}{lll}
$\mathcal T_1$ & $=(\mathrm{fpqc})=$ & \Term{faithfully-flat}{忠实平坦}\Term{quasi-compact}{拟紧}\Term{topology}{拓扑}，\\
$\mathcal T_2$ & $=(\mathrm{fppf})=$ & \Term{faithfully-flat}{忠实平坦}（\Term{local-scheme}{局部}）\Term{finite-presentation}{有限表示}\Term{topology}{拓扑}，\\
$\mathcal T_3$ & $=(\text{ét})=$ & \Term{etale}{平展}\Term{topology}{拓扑}，\\
$\mathcal T_4$ & $=(\text{étf})=$ & \Term{finite}{有限}\Term{etale}{平展}\Term{topology}{拓扑}。
\end{tabular}
\end{center}
因为 $P'_1\supset P'_2\supset P'_3\supset P'_4$，有
\[
(\mathrm{fpqc})\geq(\mathrm{fppf})\geq(\text{ét})\geq(\text{étf})\geq(\mathrm{Zar}).
\]
\begin{SGAresult}{命题}{6.3.1}
\begin{enumerate}
\item[(i)] $S$ 的\Term{sieve}{筛} $R$ 对 $\mathcal T_i$（$1\leq i\leq3$）为\Term{covering-sieve}{覆盖筛}，当且仅当存在 $S$ 的一个\Term{affine-open}{仿射开集}\Term{covering}{覆盖} $(S_p)$，以及对每个 $p$ 一个属于 $P'_i$ 的族 $\{S_{pq}\to S_p\}$，其中 $S_{pq}$ \Term{affine}{仿射}，使每个\Term{morphism}{态射} $\{S_{pq}\to S\}$ 都经 $R$ \Term{resolution}{分解}。\footnote{编者注：由假设，每个族 $\{S_{pq}\to S\}\in P'_i$ 都是\Term{finite}{有限}的，因而 $S'_p=\coprod_q S_{pq}$ \Term{affine}{仿射}，故可用仍属于 $P'_i$ 的\Term{morphism}{态射} $S'_p\to S_p$ 代替此族。}
\item[(ii)] $(\mathrm{Sch})$ 上的\Term{presheaf}{预层} $F$ 为(fpqc)（相应地，(fppf)、(ét)、(étf)）的\Term{sheaf}{层}，当且仅当：
\begin{enumerate}
\item[a)] $F$ 为(Zar)的\Term{sheaf}{层}，即一个\Term{local-nature-functor}{局部性函子}。
\item[b)] 对每个\Term{faithfully-flat}{忠实平坦}的（相应地，\Term{faithfully-flat}{忠实平坦}、\Term{finite-presentation}{有限表示}且\Term{quasi-finite}{拟有限}的，\Term{etale}{平展}\Term{surjectivity}{满射}的，\Term{finite}{有限}\Term{etale}{平展}\Term{surjectivity}{满射}的）\Term{morphism}{态射} $T\to S$，其中 $T$ 与 $S$ 均\Term{affine}{仿射}，有\Term{exact}{正合}图式
\[
F(S)\longrightarrow F(T)\rightrightarrows F(T\times_S T).
\]
\end{enumerate}
\item[(iii)] \Term{topology}{拓扑} $\mathcal T_i$（$i=1,2,3,4$）不比\Term{canonical-topology}{典范拓扑}更细。
\item[(iv)] 每个由\Term{flat}{平坦}且开的（相应地，\Term{flat}{平坦}且\Term{locally-finite-presentation}{局部有限表示}的，\Term{etale}{平展}的，\Term{finite}{有限}\Term{etale}{平展}的）\Term{morphism}{态射}构成的\Term{surjective-family}{满射族}，均为(fpqc)（相应地，(fppf)、(ét)、(étf)）的\Term{covering-family}{覆盖族}。
\SourcePage{248}
\item[(v)] 每个由\Term{flat}{平坦}且\Term{quasi-compact}{拟紧}的\Term{morphism}{态射}构成的\Term{finite}{有限}\Term{surjective-family}{满射族}，均为(fpqc)的\Term{covering-family}{覆盖族}。特别地，每个\Term{faithfully-flat}{忠实平坦}且\Term{quasi-compact}{拟紧}的\Term{morphism}{态射}都是(fpqc)的\Term{covering-morphism}{覆盖态射}。
\end{enumerate}
\end{SGAresult}
\begin{proof}
(i)由6.2.1得出，(ii)由6.2.3得出，这里要注意\Term{zariski-topology}{Zariski拓扑}的\Term{sheaf}{层}将\Term{direct-sum}{直和}变为\Term{product}{积}。每个\Term{representable}{可表}\Term{functor}{函子}都是(Zar)的\Term{sheaf}{层}，且由 SGA 1, VIII 5.3 满足(ii)的条件(b)，故 $\mathcal T_1$ 不比\Term{canonical-topology}{典范拓扑}更细，这就证明了(iii)。

证明(iv)。设 $\{S_i\to S\}$ 为命题所述的\Term{morphism}{态射}族。考虑 $S$ 的\Term{affine-open}{仿射开集}\Term{covering}{覆盖}，立即归结到 $S$ \Term{affine}{仿射}的情形。\footnote{编者注：利用现在已假设 $S$ \Term{affine}{仿射}这一点，我们简化了后面的论证。}

先处理\Term{morphism}{态射} $S_i\to S$ \Term{flat}{平坦}且开（相应地，\Term{etale}{平展}）的情形。取 $S_i$ 的\Term{affine-open}{仿射开集}\Term{covering}{覆盖} $S_{ij}$。因为所涉\Term{morphism}{态射}都是开\Term{morphism}{态射}，$S_{ij}$ 的像 $T_{ij}$ 构成 $S$ 的一个\Term{open-covering}{开覆盖}。因为 $S$ \Term{affine}{仿射}，故\Term{quasi-compact}{拟紧}，它由\Term{finite}{有限}多个\Term{open-subset}{开集} $T_{ij}$ \Term{covering}{覆盖}，其中 $(i,j)$ 遍历一个\Term{finite}{有限}集 $F$。于是 $S'=\coprod_F S_{ij}$ \Term{affine}{仿射}，\Term{morphism}{态射} $S'\to S$ 属于 $P'_1$，相应地属于 $P'_3$，因而是\Term{covering-morphism}{覆盖态射}。它经给定族\Term{resolution}{分解}，故给定族是\Term{covering-family}{覆盖族}。

在(étf)的情形，每个 $S_i$ 在 $S$ 上\Term{finite}{有限}，因而\Term{affine}{仿射}；在上一论证中，便可取 $\{S_i\}$ 为 $S_i$ 的\Term{covering}{覆盖}，并得到 $S'\to S$ 属于 $P'_4$。

现在考虑\Term{morphism}{态射} $f_i:S_i\to S$ \Term{flat}{平坦}且\Term{locally-finite-presentation}{局部有限表示}的情形。对每个 $s\in S$，由 EGA IV$_4$, 17.16.2（的证明），某个 $S_i$ 中存在\Term{affine}{仿射}\Term{subscheme}{子概形} $X(s)$，使 $s\in f_i(X(s))$，且 $f_i$ 的\Term{restriction}{限制}\Term{morphism}{态射} $g_i:X(s)\to S$ \Term{flat}{平坦}、\Term{finite-presentation}{有限表示}且\Term{quasi-finite}{拟有限}。于是，$g_i(X(s))$ 是 $s$ 的一个\Term{open-neighborhood}{开邻域} $U(s)$（EGA IV$_2$, 2.4.6），而 $S$ \Term{affine}{仿射}，因而由\Term{finite}{有限}多个这样的\Term{open-subset}{开集} $U(s_j)$（$j=1,\ldots,n$）\Term{covering}{覆盖}。因此 $X'=\coprod_j X(s_j)$ \Term{affine}{仿射}，\Term{morphism}{态射} $X'\to S$ 为\Term{surjectivity}{满射}、\Term{flat}{平坦}、\Term{finite-presentation}{有限表示}且\Term{quasi-finite}{拟有限}，故属于 $P'_2$。\footnote{编者注：这说明，若 $P''_2$ 表示由 $\mathcal C'$ 中\Term{finite-presentation}{有限表示}的\Term{flat}{平坦}\Term{morphism}{态射}构成的\Term{finite}{有限}\Term{surjective-family}{满射族}的\Term{set}{集合}，则 $P$ 与 $P''_2$ 生成的\Term{topology}{拓扑}等于 $\mathcal T_2$。另一方面，沿用(iv)证明开头的记号，若取 $S_i$ 的一个由在 $S$ 上\Term{finite-presentation}{有限表示}的\Term{affine-open}{仿射开集}构成的\Term{covering}{覆盖}，就得到 $S'\to S$ 属于 $P''_2$。}这就证明了(iv)。

证明(v)。设 $\{S_i\to S\}$ 为\Term{finite}{有限}的\Term{faithfully-flat}{忠实平坦}且\Term{quasi-compact}{拟紧}的族。取 $S$ 的\Term{affine-open}{仿射开集}\Term{covering}{覆盖} $T_j$。各 $S_{ij}=T_j\times_S S_i$ \Term{quasi-compact}{拟紧}，所以具有\Term{finite}{有限}的\Term{affine-open}{仿射开集}\Term{covering}{覆盖} $T_{ijk}$。每个\Term{morphism}{态射} $T_{ijk}\to T_j$ 都\Term{flat}{平坦}，且族 $\{T_{ijk}\to T_j\}$ \Term{finite}{有限}且\Term{surjectivity}{满射}，故是 $\mathcal T_1$ 的\Term{covering-family}{覆盖族}。于是由复合，族 $\{T_{ijk}\to S\}$ 也是\Term{covering-family}{覆盖族}。它经给定族\Term{resolution}{分解}，故给定族也是\Term{covering-family}{覆盖族}：
\SourcePage{249}
\[
\begin{tikzcd}
S_i\arrow[d]&S_{ij}\arrow[l]\arrow[d]&T_{ijk}\arrow[l]\arrow[dl]\\
S&T_j\arrow[l]&.
\end{tikzcd}
\]
\end{proof}
\begin{SGAresult}{推论}{6.3.2}
设 $(\mathrm M_i)$ 为如下\Term{morphism}{态射}族：
\begin{itemize}
\item $(\mathrm M_1)$：\Term{faithfully-flat}{忠实平坦}且\Term{quasi-compact}{拟紧}的\Term{morphism}{态射}。
\item $(\mathrm M_2)$：\Term{faithfully-flat}{忠实平坦}且\Term{locally-finite-presentation}{局部有限表示}的\Term{morphism}{态射}。
\item $(\mathrm M_3)$：\Term{etale}{平展}\Term{surjectivity}{满射}\Term{morphism}{态射}。
\item $(\mathrm M_4)$：\Term{finite}{有限}\Term{etale}{平展}\Term{surjectivity}{满射}\Term{morphism}{态射}。\footnote{编者注：我们修正了原文，为 $(\mathrm M_3)$ 和 $(\mathrm M_4)$ 加上\Term{surjectivity}{满射}假设；在其他情形下，这一假设自动成立。}
\end{itemize}
族 $(\mathrm M_i)$ 满足4.6.3的公理(a)、(b)、(c)、$(d_{\mathcal T_i})$与$(e_{\mathcal T_i})$。
\end{SGAresult}
事实上，对(a)、(b)、(c)，这是经典结果（参见 EGA 与 SGA 1 各处）。\footnote{编者注：关于“\Term{quasi-compact}{拟紧}”，参见 EGA I, 6.6.4；关于“\Term{finite}{有限}”，参见 EGA II, 6.1.5；关于“\Term{locally-finite-presentation}{局部有限表示}”，参见 EGA IV$_1$, 1.6.2；关于“\Term{faithfully-flat}{忠实平坦}”，参见 EGA IV$_2$, 2.2.13；关于“\Term{etale}{平展}”，参见 EGA IV$_4$, 17.3.3。}由6.3.1(iv)与(v)，$(\mathrm M_i)$ 满足$(d_{\mathcal T_i})$。还须说明 $(\mathrm M_i)$ 满足$(e_{\mathcal T_i})$；为此，只须说明 $(\mathrm M_i)$ 满足$(e_{\mathcal T_1})$，后者蕴含其他条件。这由 SGA 1, VIII（第4、5节）得出。
\begin{SGAresult}{推论}{6.3.3}
若 $X$ 为\Term{scheme}{概形}，$R$ 为 $X$ 中的一个 $(\mathrm M_i)$ 型\Term{equivalence-relation}{等价关系}，则 $R$ 为\Term{M-effective}{$(\mathrm M_i)$-有效的}，当且仅当 $X$ 关于 $R$ 的 $\mathcal T_i$-\Term{quotient-sheaf}{商层}\Term{representable}{可表}；此时它由\Term{quotient}{商} $X/R$ 表示。
\end{SGAresult}
这正是4.6.5。
\SGAsection{6.4}{有效性条件}
现在寻找满足4.7的公理$(f_{\mathcal T})$的\Term{morphism}{态射}族 $(\mathrm N)$。首先注意，$(f_{\mathcal T_1})$ 蕴含$(f_{\mathcal T_i})$，故可以只考虑(fpqc)\Term{topology}{拓扑}。
\begin{SGAresult}{\Term{lemma}{引理}}{6.4.1}
下列\Term{morphism}{态射}族满足4.7的公理$(f_{\mathcal T_1})$，即“可通过(fpqc)\Term{descent}{下降}”：
\begin{itemize}
\item $(\mathrm N)$：\Term{open-immersion}{开浸入}。
\item $(\mathrm N')$：\Term{closed-immersion}{闭浸入}。
\item $(\mathrm N'')$：\Term{quasi-compact}{拟紧}\Term{immersion}{浸入}。
\end{itemize}
\end{SGAresult}
\SourcePage{250}
由6.3.1(ii)，只须验证给定各族可通过\Term{zariski-topology}{Zariski拓扑}及一个\Term{faithfully-flat}{忠实平坦}\Term{quasi-compact}{拟紧}\Term{morphism}{态射}\Term{descent}{下降}。第一点显然，下面验证第二点。对 $(\mathrm N)$，这是 SGA 1, VIII 4.4；对 $(\mathrm N')$，这是同处1.9。对 $(\mathrm N'')$，利用前两个结果，按同处5.5的方式论证。
\begin{SGAresult}{推论}{6.4.2}
同样的结果对\Term{quasi-compact}{拟紧}\Term{open-immersion}{开浸入}成立。
\end{SGAresult}
由这些结果，可以在当前情形中应用4.7.1、4.7.2、5.1.8、5.3.1等处的一般结果。下面以第一个结果为例加以陈述。
\begin{SGAresult}{推论}{6.4.3}
（$=4.7.1+4.6.10$。）设 $X$ 为\Term{scheme}{概形}，$R$ 为 $X$ 中的\Term{equivalence-relation}{等价关系}。假设 $R\to X$ \Term{faithfully-flat}{忠实平坦}且\Term{quasi-compact}{拟紧}，$R\to X\times X$ 为\Term{closed-immersion}{闭浸入}（相应地，为\Term{open-immersion}{开浸入}、\Term{quasi-compact}{拟紧}\Term{immersion}{浸入}、\Term{quasi-compact}{拟紧}\Term{open-immersion}{开浸入}）。则(fpqc)\Term{topology}{拓扑}与\Term{canonical-topology}{典范拓扑}的\Term{quotient-sheaf}{商层} $X/R$ 相同，并且对每个\Term{scheme}{概形} $S$，有
\begin{center}
\( (X/R)(S)=\left\{\rule{0pt}{5.8em}\right. \)
\begin{minipage}[c]{0.68\linewidth}
$X_S$ 的\Term{closed-subscheme}{闭子概形}（相应地，\Term{open-subscheme}{开子概形}、\Term{retrocompact}{逆紧}\Term{subscheme}{子概形}、\Term{retrocompact}{逆紧}\Term{open-subscheme}{开子概形}）$Z$\footnotemark，在 $R\times S$ 下\Term{stable}{稳定}，使得 $Z\to S$ \Term{faithfully-flat}{忠实平坦}且\Term{quasi-compact}{拟紧}，并且图式 $R_Z\rightrightarrows Z\to S$ \Term{exact}{正合}
\end{minipage}
\( \left.\rule{0pt}{5.8em}\right\}. \)
\end{center}
\footnotetext{编者注：回顾，\Term{scheme}{概形} $T$ 的\Term{subscheme}{子概形} $Z$ 称为\Term{retrocompact}{逆紧的}，如果\Term{immersion}{浸入} $Z\hookrightarrow T$ \Term{quasi-compact}{拟紧}，参见 EGA $0_{\mathrm{III}}$, 9.1.1。}
\end{SGAresult}
\SGAsection{6.5}{主齐性纤维丛}
只指出如下术语：
\begin{center}
\begin{tabular}{ll}
\Term{topology}{拓扑}&\Term{principal-homogeneous-bundle}{主齐性纤维丛}\\\hline
(fpqc)&\Term{principal-homogeneous-bundle}{主齐性纤维丛}（不加限定语）\\
(ét)&\Term{quasi-isotrivial}{拟等平凡}的\Term{principal-homogeneous-bundle}{主齐性纤维丛}\\
(étf)&\Term{locally-isotrivial}{局部等平凡}的\Term{principal-homogeneous-bundle}{主齐性纤维丛}\\
(Zar)&\Term{locally-trivial}{局部平凡}的\Term{principal-homogeneous-bundle}{主齐性纤维丛}。
\end{tabular}
\end{center}
\SGAsection{6.6}{其他拓扑}
有时还会使用\Term{scheme}{概形}\Term{category}{范畴}上的其他\Term{topology}{拓扑}。这里指出一种：\Term{global-finite-etale-topology}{全局有限平展拓扑}(étfg)，它由如下\Term{pretopology}{预拓扑}生成：其\Term{covering-family}{覆盖族}为由\Term{finite}{有限}\Term{etale}{平展}\Term{morphism}{态射}构成的\Term{surjective-family}{满射族}。它不比\Term{zariski-topology}{Zariski拓扑}更细。相应的\Term{principal-homogeneous-bundle}{主齐性纤维丛}称为“\Term{isotrivial}{等平凡的}”。
\SourcePage{251}
\begin{center}
\begin{tikzpicture}[x=1.6cm,y=0.9cm]
\node (canonical) at (0,6) {（典范）};
\node (fpqc) at (0,5) {(fpqc)};
\node (fppf) at (0,4) {(fppf)};
\node (et) at (0,3) {(ét)};
\node (etf) at (0,2) {(étf)};
\node (zar) at (-1,1) {(Zar)};
\node (etfg) at (1,1) {(étfg)};
\node (chaotic) at (0,0) {（混沌）};
\draw (canonical)--(fpqc)--(fppf)--(et)--(etf);
\draw (etf)--(zar)--(chaotic)--(etfg)--(etf);
\end{tikzpicture}\footnote{编者注：回顾（参见4.4.2），\Term{chaotic-topology}{混沌拓扑}是最粗的\Term{topology}{拓扑}，由对所有 $S\in\operatorname{Ob}\mathcal C$ 取 $J(S)=\{S\}$ 定义。}
\end{center}
\SGAsection{6.7}{齐性空间}
\footnote{编者注：我们增补了以下各小节。}设 $G$ 为 $S$-\Term{group-scheme}{群概形}，$X$ 为具有 $G$ 的\Term{left-action}{左作用}的 $S$-\Term{scheme}{概形}，设
\[
\Phi:G\times_S X\longrightarrow X\times_S X
\]
为在\Term{set-sheaf}{集合层}面由 $(g,x)\mapsto(gx,x)$ 定义的 $S$-\Term{scheme}{概形}\Term{morphism}{态射}。回顾（参见5.1.0及III 0.1），若以下\Term{equivalence}{等价}条件成立，就称 $X$ 为在 $G$ 下的\Term{formally-principal-homogeneous-space}{形式主齐性空间}：
\begin{enumerate}
\item[(i)] 对每个 $T\to S$，\Term{set}{集合} $X(T)$ 为空或在 $G(T)$ 下\Term{principal-homogeneous}{主齐性}；
\item[(ii)] $\Phi$ 为 $S$-\Term{functor}{函子}的\Term{isomorphism}{同构}；
\item[(iii)] $\Phi$ 为 $S$-\Term{scheme}{概形}的\Term{isomorphism}{同构}。
\end{enumerate}
（(i)$\Leftrightarrow$(ii)显然；因 $\mathcal C=(\mathrm{Sch}/S)$ 是 $\widehat{\mathcal C}$ 的一个\Term{full-subcategory}{满子范畴}，故(ii)$\Leftrightarrow$(iii)。）

\Term{formally-homogeneous}{形式齐性}（未必\Term{principal-homogeneous}{主齐性}）\Term{homogeneous-space}{空间}的定义，是要求 $\Phi$ 在适当\Term{topology}{拓扑} $\mathcal T$ 的\Term{sheaf}{层}\Term{category}{范畴}中为\Term{epimorphism}{满态射}。事实上，$\Phi$ 为 $S$-\Term{functor}{函子}的\Term{epimorphism}{满态射}，\Term{equivalence}{等价}于对每个 $T\to S$，\Term{set}{集合} $X(T)$ 为空或在 $G(T)$ 下\Term{homogeneous}{齐性}（未必\Term{principal-homogeneous}{主齐性}）；但这一条件过强，下面的简单例子即可说明。设 $S=\operatorname{Spec}\mathbb R$，$G=\mathbb G_{m,\mathbb R}$，$X=\mathbb G_{m,\mathbb R}$，$G$ 通过 $t\cdot x=t^2x$ 在 $X$ 上\Term{action}{作用}。
\SourcePage{252}
则\Term{morphism}{态射} $\Phi$ \Term{etale}{平展}、\Term{finite}{有限}且\Term{surjectivity}{满射}，因而是(étf)\Term{topology}{拓扑}的\Term{sheaf}{层}\Term{category}{范畴}中的\Term{epimorphism}{满态射}（从而更是 $S$-\Term{scheme}{概形}的\Term{epimorphism}{满态射}）；但 $X(\mathbb R)$ 中的点 $1$ 与 $-1$ 不由 $G(\mathbb R)$ 的任何元素相互变换，因此\Term{morphism}{态射} $G(\mathbb R)\times X(\mathbb R)\to X(\mathbb R)\times X(\mathbb R)$ 不是\Term{surjectivity}{满射}\footnote{编者注：显然，这一困难源于如下事实：若 $\mathcal C'$ 是 $\widehat{\mathcal C}$ 的一个包含 $\mathcal C$ 的\Term{full-subcategory}{满子范畴}，例如 $\mathcal C$ 上某个不比\Term{canonical-topology}{典范拓扑}更细的\Term{topology}{拓扑} $\mathcal T$ 的\Term{sheaf}{层}\Term{category}{范畴} $\widetilde{\mathcal C}_{\mathcal T}$，且 $f:X\to Y$ 为 $\mathcal C$ 中的\Term{morphism}{态射}，则下列蕴含一般是严格的：
\[
f\text{ 是 }\widehat{\mathcal C}\text{ 的满态射}\ \Rightarrow\
f\text{ 是 }\mathcal C'\text{ 的满态射}\ \Rightarrow\
f\text{ 是 }\mathcal C\text{ 的满态射}.
\]
}。因此作如下定义：
\begin{SGAresult}{定义}{6.7.1}
设 $G$ 为 $S$-\Term{group}{群}，$X$ 为具有 $G$ 的\Term{action}{作用}的 $S$-\Term{scheme}{概形}，$\mathcal T$ 为 $(\mathrm{Sch}/S)$ 上不比\Term{canonical-topology}{典范拓扑}更细的\Term{topology}{拓扑}。若以下\Term{equivalence}{等价}条件成立，就称 $X$ 为（相对于\Term{topology}{拓扑} $\mathcal T$）在 $G$ 下的\Term{formally-homogeneous-space}{形式齐性空间}：
\begin{enumerate}
\item[(i)] \Term{morphism}{态射} $\Phi:G\times_S X\to X\times_S X$ 是\Term{topology}{拓扑} $\mathcal T$ 的\Term{sheaf}{层}\Term{category}{范畴}中的\Term{epimorphism}{满态射}；
\item[(ii)] 对每个 $T\to S$ 及 $x,y\in X(T)$，存在对\Term{topology}{拓扑} $\mathcal T$ 为\Term{covering-morphism}{覆盖态射}的\Term{morphism}{态射} $T'\to T$，以及 $g\in G(T')$，使 $y_{T'}=g\cdot x_{T'}$。
\end{enumerate}
\end{SGAresult}
\begin{SGAresult}{注记}{6.7.2}
特别地，条件(i)蕴含 $\Phi$ 为 $(\mathrm{Sch}/S)$ 中的\Term{universal-effective-epimorphism}{泛有效满态射}（参见4.4.3）。容易看出，这蕴含 $\Phi$ 为\Term{surjectivity}{满射}（参见1.3，编者注(3)）。
\end{SGAresult}
\begin{SGAresult}{命题与定义}{6.7.3}
\footnote{编者注：参见[Ray70]，定义VI.1.1。}设 $G$ 为 $S$-\Term{group}{群}，$X$ 为具有 $G$ 的\Term{action}{作用}的 $S$-\Term{scheme}{概形}，$\mathcal T$ 为 $(\mathrm{Sch}/S)$ 上不比\Term{canonical-topology}{典范拓扑}更细的\Term{topology}{拓扑}。以下条件\Term{equivalence}{等价}：
\begin{enumerate}
\item[(i)] $X$ 满足下面两个假设：
\begin{enumerate}
\item[(1)] \Term{morphism}{态射} $\Phi:G\times_S X\to X\times_S X$ 为\Term{covering-morphism}{覆盖态射}，即 $X$ 是一个\Term{formally-homogeneous}{形式齐性} $G$-\Term{homogeneous-space}{空间}；
\item[(2)] \Term{morphism}{态射} $X\to S$ 也为\Term{covering-morphism}{覆盖态射}，即相对于\Term{topology}{拓扑} $\mathcal T$，它\Term{local-scheme}{局部}具有一个\Term{section}{截面}（参见 RefIV.4.4.8bis.bis）。
\end{enumerate}
\item[(ii)] “在 $S$ 上相对于\Term{topology}{拓扑} $\mathcal T$ \Term{local-scheme}{局部}地”，$X$ 作为具有\Term{group}{群} $G$ 的\Term{action}{作用}的\Term{scheme}{概形}，\Term{isomorphism}{同构}于 $G$ 关于某个\Term{subgroup-scheme}{子群概形} $H$ 的（$\mathcal T$-）\Term{quotient-sheaf}{商层}；即存在\Term{covering-family}{覆盖族} $\{S_i\to S\}$，使每个 $X\times_S S_i$ 都表示 $G\times_S S_i$ 关于某个\Term{subgroup-scheme}{子群概形} $H_i$ 的\Term{quotient-sheaf}{商层}。
\end{enumerate}
在这些条件下，称 $X$ 为一个（相对于\Term{topology}{拓扑} $\mathcal T$ 的）$G$-\Term{homogeneous-space}{齐性空间}。
\end{SGAresult}
\begin{proof}
假设(ii)成立。令 $G_i=G\times_S S_i$，$X_i=X\times_S S_i$。则 $X_i$ 在 $S_i$ 上具有一个\Term{section}{截面}，即\Term{unit-section}{单位截面} $\varepsilon_i:S_i\to G_i$ 与\Term{projection}{投影} $\pi_i:G_i\to X_i=G_i/H_i$ 的复合。因此 $X\to S$ 为\Term{covering-morphism}{覆盖态射}。
\SourcePage{253}
另一方面，$\pi_i$ 为\Term{covering-morphism}{覆盖态射}，所以 $\pi_i\times\pi_i$ 也为\Term{covering-morphism}{覆盖态射}（参见4.2.3(C1)与(C2)），并有\Term{commutative-diagram}{交换图}：
\[
\begin{tikzcd}[column sep=large]
G_i\times_{S_i}X_i\arrow[r,"\Phi_i"]&G_i\times_{S_i}X_i\\
G_i\times_{S_i}G_i\arrow[u,"\mathrm{id}\times\pi_i"]\arrow[r,"\Psi_i","\sim"']&G_i\times_{S_i}G_i\arrow[u,"\pi_i\times\pi_i"']
\end{tikzcd}
\]
其中 $\Phi_i$ 由 $\Phi$ 经\Term{base-change}{基变换} $S_i\to S$ 得到，而 $\Psi_i$ 是在\Term{set-sheaf}{集合层}面由 $(g,g')\mapsto(gg',g)$ 定义的\Term{isomorphism}{同构}。于是 $(\pi_i\times\pi_i)\circ\Psi_i$ 为\Term{covering-morphism}{覆盖态射}，故 $\Phi_i$ 也为\Term{covering-morphism}{覆盖态射}（4.2.3(C3)）。这表明 $\Phi$“\Term{local-scheme}{局部}为\Term{covering-morphism}{覆盖态射}”，因而是\Term{covering-morphism}{覆盖态射}（4.2.3(C5)）。这就证明了(ii)$\Rightarrow$(i)。

反过来，假设(i)成立，并进一步假设\Term{structure-morphism}{结构态射} $X\to S$ 具有一个\Term{section}{截面} $\sigma$。由 EGA I, 5.3.13，$\sigma$ 是一个\Term{immersion}{浸入}。用下图定义 $H=G\times_X S$，其中两个方块均为\Term{cartesian-square}{笛卡儿方块}：
\[
\begin{tikzcd}[column sep=large]
H\arrow[r]\arrow[d]&G\arrow[r,"\mathrm{id}_G\boxtimes\sigma"]\arrow[d,"\pi"]&G\times_S X\arrow[d,"\Phi"]\\
S\arrow[r,"\sigma"']&X\arrow[r,"\mathrm{id}_X\boxtimes\sigma"']&X\times_S X
\end{tikzcd}
\]
其中 $\pi$、$\mathrm{id}_G\boxtimes\sigma$ 与 $\mathrm{id}_X\boxtimes\sigma$ 表示在\Term{set-sheaf}{集合层}面由下式定义的\Term{morphism}{态射}：对 $T\to S$ 及 $g\in G(T)$、$x\in X(T)$，
\[
\pi(g)=g\cdot\sigma_T,\qquad
(\mathrm{id}_G\boxtimes\sigma)(g)=(g,\sigma_T),\qquad
(\mathrm{id}_X\boxtimes\sigma)(x)=(x,\sigma_T).
\]
于是 $\pi$ 为\Term{covering-morphism}{覆盖态射}，$H$ 是 $G$ 的一个\Term{subgroup-scheme}{子群概形}，表示 $\sigma$ 的\Term{stabilizer}{稳定子} $\underline{\operatorname{Stab}}_G(\sigma)$（参见I 2.3.3），即对每个 $T\to S$，有
\[
H(T)=\{g\in G(T)\mid g\cdot\sigma_T=\sigma_T\}.
\]
以 $G/H$ 表示\Term{presheaf}{预层} $T\mapsto G(T)/H(T)$，以 $a(G/H)$ 表示相对于\Term{topology}{拓扑} $\mathcal T$ 的\Term{associated-sheaf}{伴随层}。由上文，得到具有\Term{group}{群} $G$ 的\Term{action}{作用}的\Term{presheaf}{预层}的\Term{morphism}{态射}所成的\Term{commutative-diagram}{交换图}：
\[
\begin{tikzcd}
G\arrow[r,"\pi"]\arrow[d]&X\\
G/H\arrow[ur,"\bar\pi"']&,
\end{tikzcd}
\]
其中 $\bar\pi$ 为\Term{monomorphism}{单态射}。因为 $\pi$ 为\Term{covering-morphism}{覆盖态射}，$\bar\pi$ 也为\Term{covering-morphism}{覆盖态射}，故由4.3.12，$\bar\pi$ 诱导\Term{isomorphism}{同构} $a(G/H)\xrightarrow{\sim}X$。因此已证明：若 $X$ 是一个 $G$-\Term{homogeneous-space}{齐性空间}，且 $X\to S$ 有一个\Term{section}{截面} $\sigma$，则 $X$ 表示\Term{quotient-sheaf}{商层} $G/H$，其中 $H=G\times_X S$ 为 $\sigma$ 的\Term{stabilizer}{稳定子}。

在一般情形，依假设存在\Term{covering-family}{覆盖族} $\{S_i\to S\}$，使每个\Term{morphism}{态射} $X_i=X\times_S S_i\to S_i$ 具有一个\Term{section}{截面} $\sigma_i$。令 $G_i=G\times_S S_i$；则由 $\Phi$ 经\Term{base-change}{基变换}
\SourcePage{254}
$S_i\to S$ 得到的\Term{morphism}{态射} $\Phi_i:G_i\times_{S_i}X_i\to X_i\times_{S_i}X_i$ 仍为\Term{covering-morphism}{覆盖态射}。因此，由上文有 $X_i\simeq G_i/H_i$，其中 $H_i$ 是 $\sigma_i$ 在 $G_i$ 中的\Term{stabilizer}{稳定子}。这就完成了(i)$\Rightarrow$(ii)的证明。
\end{proof}
\section*{参考文献}\addcontentsline{toc}{section}{参考文献}
\begin{description}
\item[{[AS]}] J. Giraud，\textit{Analysis Situs}（《位置分析》），Sém. Bourbaki，报告256，1963年5月。
\item[{[D]}] J. Giraud，\textit{Méthode de la descente}（《\Term{descent}{下降}方法》），Mém. Soc. Math. France，第2卷（1964），第iii--viii及1--150页。
\item[{[MA]}] M. Artin，\textit{Grothendieck Topologies}（《Grothendieck\Term{topology}{拓扑}》），油印讲义，Harvard，1962。
\item[{[SGA 1]}] \textit{Séminaire de Géométrie Algébrique du Bois-Marie 1960--61, Revêtements étales et groupe fondamental}（《Bois-Marie\Term{algebraic-geometry}{代数几何}讨论班1960--61：\Term{etale}{平展}\Term{covering}{覆盖}与\Term{fundamental-group}{基本群}》），Lecture Notes in Maths. 224（1971）；重排并加注版，Documents Math. 3，Soc. Math. France，2003。
\item[{[SGA 4]}] \textit{Séminaire de Géométrie Algébrique du Bois-Marie 1963--1964, Théorie des topos et cohomologie étale des schémas}（《Bois-Marie\Term{algebraic-geometry}{代数几何}讨论班1963--1964：\Term{topos}{拓扑斯}理论与\Term{scheme}{概形}的\Term{etale}{平展}\Term{cohomology}{上同调}》），第I、II、III卷，Lecture Notes in Maths. 269、270（1972），305（1973）。
\item[{[TDTE I]}] A. Grothendieck，\textit{Techniques de descente et théorèmes d'existence en géométrie algébrique I. Généralités. Descente par morphismes fidèlement plats}（《\Term{algebraic-geometry}{代数几何}中的\Term{descent}{下降}技术与\Term{existence-theorem}{存在定理}I：一般理论；沿\Term{faithfully-flat}{忠实平坦}\Term{morphism}{态射}的\Term{descent}{下降}》），Sém. Bourbaki，报告190，1959年12月。
\item[{[Ray70]}] \footnote{编者注：我们补入了这一文献。}M. Raynaud，\textit{Faisceaux amples sur les schémas en groupes et les espaces homogènes}（《\Term{group-scheme}{群概形}与\Term{homogeneous-space}{齐性空间}上的\Term{ample}{丰沛}\Term{sheaf}{层}》），Lect. Notes Math. 119，Springer-Verlag，1970。
\end{description}
